\documentclass[11pt,a4paper]{article}
\usepackage[T1]{fontenc}
\usepackage[utf8]{inputenc}
\usepackage{lmodern,amsmath,amssymb,array,longtable,microtype}
\usepackage[margin=24mm]{geometry}
\usepackage[colorlinks=true,urlcolor=blue,linkcolor=blue]{hyperref}
\usepackage{xurl}
\setlength{\parindent}{0pt}
\setlength{\parskip}{6pt}
\setlength{\emergencystretch}{2em}
\newcommand{\code}[1]{\texttt{\detokenize{#1}}}
\title{Preparing the Formalization\\\large Analytic interfaces, calculations, and reusable resources}
\author{}
\date{September 17, 2026}
\begin{document}
\maketitle
\begin{quote}\small
Historical preparatory note. This English translation preserves the original
plan and its status as of September 17, 2026. For the completed formalization
and its two remaining classical admissions, see the repository README and
\path{docs/ADMISSIONS.md}.
\end{quote}
This note accompanies the expanded proofs of \emph{Double-wells in a strong
magnetic field: absence of tunneling for infinitely many values of the coupling
constant}, in the supplied 56-page PDF version. It does not change the article's
statements. Information about repositories is based on inspection of their
public files and official Mathlib documentation as of the date above; the
external repositories were not compiled here.

\section{Deliverables and navigation}
The PDF's 57 proof blocks, including appendices, are expanded into 176
sublemmas. The three results whose full proofs are in appendices retain that
placement: the main-text block explains the assembly and refers to the appendix
sublemmas. These are therefore not three unfinished proofs.

Each new statement actually uses \verb|\begin{lemma}| and a label beginning
with \code{sublemma:}. A local environment displays it as a sublemma, with
numbering independent of the article's lemmas. For example,
\path{sublemma:D6-4-phase-difference} denotes the exact phase calculation near
the saddle point; its prefix and attachment to proof 6.4 remain stable for a
future Lean translation.

The integrated PDF contains 93 pages. Since the \LaTeX{} source of that version
was not supplied, material outside proofs is preserved by vector inclusion of
fragments of the original PDF: it is neither rewritten nor replaced by an
approximate text extraction. The blue proofs are actual editable \LaTeX{} text.
The standalone proof collection has 63 pages. Figures, statements, editorial
comments, and existing colors outside proofs remain those of the supplied PDF.

Physical pagination changes. To avoid altering material outside proofs, the
page numbers visible in the original table of contents are preserved: they
still refer to pages of the initial PDF. Proof bookmarks and
\path{correspondance_preuves.csv} give the new pagination;
\path{index_sous_lemmes.csv} gives the page of each of the 176 sublemmas.
Original equation numbers are preserved in the proofs. Files
\path{proofs/*.tex} are separated by proof, and \path{preuves_a_integrer.tex}
collects them in manuscript order.

\textbf{Status.} These are natural-language proofs intended to prepare
formalization, not proofs already verified by Lean. \LaTeX{} compilation,
label uniqueness, and internal-reference resolution were checked. Those
checks do not constitute mathematical certification. Classical analytic inputs
remain identified as such.

\section{Proposed boundary: formal calculations and human interfaces}
The decomposition should not take the article's result as a disguised
hypothesis. An interface should concern a general analytic operation;
constants, signs, powers of $h$, action margins, and assemblies specific to this
construction should remain proved. The following table specifies the contracts
to isolate.

\small
\begin{longtable}{@{}>{\raggedright\arraybackslash}p{.18\textwidth}>{\raggedright\arraybackslash}p{.35\textwidth}>{\raggedright\arraybackslash}p{.39\textwidth}@{}}
\textbf{Interface}&\textbf{Contract to make explicit}&\textbf{Calculations retained in the formalized core}\\\hline
\endfirsthead
\textbf{Interface}&\textbf{Contract to make explicit}&\textbf{Calculations retained in the formalized core}\\\hline
\endhead
Harmonic approximation&For the fixed radial well: first-level asymptotics, simplicity/radiality of the ground state, and local convergence after rescaling. Check the assumptions of the result cited as [8] in the article.&Potential Taylor expansion, dilation conjugation, oscillator Gaussian, positivity of the limiting gap, and inequality transfer: proof L2.1.\\[7pt]
Landau resolvent&Self-adjoint realization and resolvent kernel at energy strictly below the spectrum; validity of the integral representation and then its radial projection.&Critical time, action value and derivatives, gauge-phase signs, change of variables, and factor $2\pi$: P2.2, P5.1, LB.2.\\[7pt]
Forms and spectrum&Closed-form representation theorem, variational characterization, spectral projections, and invertibility criterion for a coercive compression.&Rank-one decomposition, completion of the square, Schur identity, scalar monotonicity, and level comparison: Sections 3, 4, and 8.\\[7pt]
IMS and Agmon&Extension of product identities to the form domain, approximation of Lipschitz weights, and passage to the limit for saturated weights.&Exact identities, term $-\kappa^2|\nabla T|^2$, orthogonality defect, and quantitative absorption: L2.5, L3.1, LA.1--LA.2.\\[7pt]
Interior regularity&On two fixed concentric balls: an $H^{k+2}$ estimate from $L^2$ and an $H^k$ right-hand side, uniform for the bounded coefficients considered; embedding into $C^k$ in dimension two.&Local gauge, rescaling $x=x_0+hy$, weight oscillation, powers of $h$, and return to original derivatives: LA.3, LB.1.\\[7pt]
General complex analysis&Holomorphy of dominated integrals, local holomorphic inversion, Cauchy's theorem, and derivative estimates. A substantial part can connect directly to Mathlib.&Saddle-point construction, branch choices, contours, connectors, Gaussian moments, and remainder bounds: Appendices C and D, P6.7.\\
\end{longtable}
\normalsize
The interfaces do not replace oscillatory evaluation or control of the
scattered term. In particular, the weighted-inverse bound must not simply be
assumed: LA.2 gives its proof from the form domain and IMS. Likewise, the
log-flat integral's asymptotic formula is not an admitted interface: it is
proved in Appendix D.

A useful organization of the future development would put scalar and geometric
calculations in independent modules, then prove the abstract spectral-transfer
theorem in Hilbert space, and finally instantiate the magnetic problem's
analytic contracts. Conditional results should explicitly display
unformalized hypotheses; a global \code{axiom} whose presence becomes invisible
downstream should be avoided.

\section{Points made explicit in the proofs}
\textbf{Complex branches and cutoffs.} The development distinguishes variables
actually extended to the complex plane from tangential parameters that stay
real. The smooth function $\chi_a$ is not assumed holomorphic: only the formula
in the region where it equals one is extended. The scattered source is never
extended. The principal logarithm and square root are chosen on fixed
neighborhoods avoiding their cuts. For $w=W_0(z)$, the equation
$w+\operatorname{Log}w=\operatorname{Log}z$ fixes the branch before imaginary
parts are taken.

\textbf{Actual saddle-point remainder control.} The identity
\[
 f(y_c+q)-f(y_c)=\beta q^2+2\beta w(e^{-q}-1+q)
\]
is isolated. The perturbative expansion is made on a shrinking window
$|q|\le(\Re w)^{-2/5}$, not a fixed neighborhood where the cubic term could grow
large. Cancellation of the odd term and fourth and sixth moments then give
the relative remainder $O(|w|^{-1})$. The absolute contour-integral bound is
also established; it is essential for inserting the joint-phase remainder in
P6.7.

\textbf{Quantifiers and fixed potential.} Geometric parameters and the
potential are fixed before choosing $L$. After fixing $L$, one may shrink an
analytic neighborhood or the threshold $h_0(L)$, but not silently shorten the
potential's support. This distinction is explicit in L6.3 and the final proof.
Parameter-uniform estimates are always restricted to the compact sets actually
used.

\textbf{Positive scale and passage to crossings.} Errors are compared with the
positive envelope $a_h$, never with the hopping coefficient, which can vanish.
Polynomial and log-quadratic losses are absorbed only after preserving a
strictly positive margin. The $h^{-13/2}$ factor calculation, three log-flat
costs against two, and phase derivative are separated. The intermediate-value
argument assumes no differentiability of the oscillatory remainder.

\textbf{Continuity and dilations.} In L8.6, continuity of spectral bottoms is
deduced from form comparison on a fixed domain. Resolvent norm continuity is
used for the dilated family. Return to physical vectors uses strong continuity
of dilations on a fixed vector; it does not incorrectly infer operator-norm
continuity from that strong continuity alone. Rank-one projections, local
normalization, and hopping's independence of phase choice are detailed.

\textbf{Radial normalization.} LB.2 separates Sturm--Liouville factorization,
identification of the kernel constant with measure $2\pi s\,ds$, positivity of
the regular solution, and the source-integral lower bound. This prevents a
missing factor in $\Gamma_h$ from being hidden by later exponential comparisons.

\section{Directly relevant Mathlib resources}
\subsection{Real and complex Gaussians}
The module
\begin{center}\small\path{Mathlib.Analysis.SpecialFunctions.Gaussian.GaussianIntegral}\end{center}
contains \code{integral_gaussian_complex}, under $0<\Re b$, as well as
\path{integrable_cexp_neg_mul_sq}, \path{norm_cexp_neg_mul_sq}, and
\path{integrable_rpow_mul_exp_neg_mul_sq}. This is a concrete entry point for
the complex Gaussian and moment majorants in Appendix D. It also contains
asymptotic exponential-decay comparisons. Branch choice and substitution
$b=\beta(1+w)$ still need connecting to the manuscript's sublemmas. See the
\href{https://leanprover-community.github.io/mathlib4_docs/Mathlib/Analysis/SpecialFunctions/Gaussian/GaussianIntegral.html}{official module documentation}.

\subsection{Rectangular contours and Cauchy estimates}
In \path{Mathlib.Analysis.Complex.CauchyIntegral}, the entries
\begin{center}\small
\path{Complex.integral_boundary_rect_eq_zero_of_differentiableOn}\\
\path{DifferentiableOn.deriv_eq_smul_circleIntegral}
\end{center}
correspond to the rectangular contour in the logarithmic variable and
remainder-derivative estimates. The limit of the vertical face at infinity
and its majorants must be proved separately, as in D6.4. Appendix C's Morse
curves require more adaptation than this rectangle. See the
\href{https://leanprover-community.github.io/mathlib4_docs/Mathlib/Analysis/Complex/CauchyIntegral.html}{official CauchyIntegral documentation}.

\subsection{Hilbert-space projections}
The module \path{Mathlib.Analysis.InnerProductSpace.Projection.Basic}
provides, in particular, \path{Submodule.orthogonalProjectionOnto} and
\path{Submodule.starProjection}. The latter is an endomorphism of the ambient
space, especially suited to $P+Q=I$. This library can support the projection
algebra in Sections 3--4; it does not by itself provide a realization of the
unbounded magnetic Schr\"odinger operator. The documentation consulted also
marks the old name \code{orthogonalProjection} as deprecated: the name to use
therefore depends on the pinned version. See the
\href{https://leanprover-community.github.io/mathlib4_docs/Mathlib/Analysis/InnerProductSpace/Projection/Basic.html}{official projection documentation}.

\section{What to reuse from Scott Armstrong's repositories}
\subsection{CoarseGraining: Sobolev spaces, products, and Hilbert realizations}
Repository: \url{https://github.com/scottnarmstrong/CoarseGraining}.

\href{https://github.com/scottnarmstrong/CoarseGraining/blob/main/Homogenization/Sobolev/H1.lean}{\path{Homogenization/Sobolev/H1.lean}}
is an entry point, but not a standalone file: it imports the H1/Algebra
subtree. The file actually consulted,
\begin{center}\small\path{Homogenization/Sobolev/H1/Algebra/H1Function.lean}\end{center}
contains algebraic structure on real $H^1$ functions and the product rule with
a smooth compactly supported function. Entries
\path{mulContDiffHasCompactSupport} and
\path{mulContDiffHasCompactSupport_grad} are especially relevant for
localization. Reuse their imports, not only these two declarations. See the
\href{https://github.com/scottnarmstrong/CoarseGraining/blob/main/Homogenization/Sobolev/H1/Algebra/H1Function.lean}{source file}.

\href{https://github.com/scottnarmstrong/CoarseGraining/blob/main/Homogenization/Sobolev/Foundations/CoerciveH10.lean}{\path{Homogenization/Sobolev/Foundations/CoerciveH10.lean}}
exposes smooth approximations of $H^1_0$ functions and convergence in value
and gradient norms. This is a possible route for density arguments. Note that
the coercive estimates displayed in this file explicitly assume $2<d$; they
cannot be applied unchanged in dimension two.

Finally,
\href{https://github.com/scottnarmstrong/CoarseGraining/blob/main/Homogenization/PDE/HarmonicHilbert.lean}{\path{Homogenization/PDE/HarmonicHilbert.lean}}
gives a Hilbert realization of harmonic gradients, a closed subspace, a
bilinear form, and coercivity lemmas. It provides an architectural model for
passing from Sobolev objects to bounded operators between Hilbert spaces. Its
subject nevertheless remains a real divergence-form elliptic problem, and
some constructions take data such as \code{PotentialSolenoidalL2Data} as
parameters: read the fields of these data rather than treating the theorem
name alone as a sufficient contract.

\subsection{DeGiorgi: rescaling and localization}
Repository: \url{https://github.com/scottnarmstrong/DeGiorgi}.

\href{https://github.com/scottnarmstrong/DeGiorgi/blob/main/DeGiorgi/BallScaling.lean}{\path{DeGiorgi/BallScaling.lean}}
contains affine transports between a ball and the unit ball, notably
\code{rescaleToUnitBall}, \code{transportFromUnitBall}, their inverse identity,
and transport of weak formulations. This is directly relevant as a model for
$x=x_0+hy$ in LA.3. Some auxiliary measure identities are declared
\code{private}: they are not automatically available under a public name
through a simple import. The file
\href{https://github.com/scottnarmstrong/DeGiorgi/blob/main/DeGiorgi/Localization.lean}{\path{DeGiorgi/Localization.lean}}
complements this route for restriction to subdomains.

The principal limitation is precise: \code{weak_harnack_on_ball}, in
\href{https://github.com/scottnarmstrong/DeGiorgi/blob/main/DeGiorgi/ScaledBallEstimates.lean}{\path{DeGiorgi/ScaledBallEstimates.lean}},
concerns positive real functions and assumes $2<(d:\mathbb R)$. This does
not replace the $H^{k+2}\to C^k$ estimate for the two-dimensional complex
magnetic equation in LA.3. Reusing transports, function spaces, and product
rules first is more direct than trying to apply Harnack to this problem.

\subsection{Versions: two different environments}
The files consulted specify \code{leanprover/lean4:v4.29.0-rc6} for DeGiorgi,
with Mathlib pinned to commit
\begin{center}\small\path{5c8398df528176d9c87ccd9226ba8f7c8852d59c}\end{center}
CoarseGraining specifies \code{leanprover/lean4:v4.33.0}, with Mathlib at
\begin{center}\small\path{db584cd6d46c92f209a44c0f1c829460d327499d}\end{center}
These values come from DeGiorgi's
\href{https://github.com/scottnarmstrong/DeGiorgi/blob/main/lean-toolchain}{lean-toolchain}
and \href{https://github.com/scottnarmstrong/DeGiorgi/blob/main/lake-manifest.json}{lake-manifest},
and CoarseGraining's
\href{https://github.com/scottnarmstrong/CoarseGraining/blob/main/lean-toolchain}{lean-toolchain}
and \href{https://github.com/scottnarmstrong/CoarseGraining/blob/main/lake-manifest.json}{lake-manifest}.
Main branches evolve: also record the repository's own commit when retrieving it.

These trees should not be mixed without porting. Compile each repository first
with its own version files, then choose one target environment and port a
limited subset of dependencies. The inspection performed here replaces
neither that compilation nor inspection of transitive axioms and structure
hypotheses. Preserve license notices in reused files.

\section{Realistic scope of an initial partial formalization}
The first batch can cover action and geometry calculations (P2.2, P2.6,
L2.7--L2.10), phase signs and symmetries (P5.1, L5.2--L5.3), the full scalar
log-flat calculation (Appendix D), then Feshbach--Schur algebra and envelope
transfer. These components actually reduce human verification work without
immediately requiring an entire magnetic-PDE library.

In the specific modules inspected, I did not identify a ready-made package
covering the full chain ``unbounded magnetic operator -- weighted Agmon --
Landau resolvent -- harmonic approximation -- crossings.'' This search finding
is not proof of absence throughout Mathlib or external repositories. The
proposed decomposition allows calculations to be formalized without relying
on such an absence claim.

For each final conditional theorem, publish the exact list of still-external
interfaces. This distinguishes Lean verification of the calculations and their
assembly from a complete formalization of the article's theorem. Both goals
are useful, but should not be presented as equivalent.
\end{document}
